problem:  
Let \(M^n\) be a complete, noncompact Riemannian manifold with nonnegative sectional curvature and compact soul \(S \subset M\).  
Let \(\pi: M \to S\) be the Sharafutdinov retraction, and let \(F_p = \pi^{-1}(p)\) denote the totally convex fiber over \(p \in S\).  
Fix a sequence of points \(x_i \in M\) with \(r_i := \operatorname{dist}(S,x_i) \to \infty\), and define the rescaled pointed manifolds  
\[
(M_i, g_i, x_i) := (M, r_i^{-2} g, x_i).
\]
Suppose that for every fixed \(p \in S\) and every sequence \(y_i \in F_p\) with \(\operatorname{dist}(S,y_i) = r_i\), the rescaled fibers \((F_p, r_i^{-2} g|_{F_p}, y_i)\) converge in the pointed Cheeger–Gromov sense to Euclidean space \((\mathbb{R}^k, g_{\mathrm{eucl}}, 0)\), where \(k = n - \dim S\).  

**Question:**  
Does such fiberwise asymptotic Euclideanity force the blow-down limits of \((M, r_i^{-2} g, x_i)\) to be globally isometric to a Riemannian product \((S_\infty \times \mathbb{R}^k, g_\infty + g_{\mathrm{eucl}})\)?  
Equivalently, can Euclidean asymptotics along each normal fiber determine that \(M\) is asymptotically a metric product of its soul and a Euclidean factor under \(\sec M \ge 0\)?  

Why is it a "good" problem:  
This formulation eliminates redundant assumptions and reduces to a core rigidity question: does **local Euclideanity of Sharafutdinov fibers** imply **global product structure** in the asymptotic geometry of nonnegatively curved manifolds?  

It is good for several reasons:  

1. **Conceptual Simplicity, Profound Depth:**  
   The statement is simple—if fibers look Euclidean at infinity, must the whole manifold look like a product?—yet addressing it requires deep analysis of how sectional curvature couples horizontal and vertical directions via the **Riccati equations** governing the shape of distance tubes around \(S\).  

2. **Bridges Distinct Realms:**  
   The problem connects **comparison geometry** (curvature inequalities and splitting theorems), **metric convergence theory** (Cheeger–Gromov blow-down limits), and **submersion geometry** (behavior of Sharafutdinov retraction fibers and their parallel transport). Any progress would deepen understanding of how fiber flatness propagates horizontally—a phenomenon presently understood only in special cases (e.g., totally flat or symmetric situations).  

3. **Potential for a New Rigidity Principle:**  
   If true, the result would extend the **Cheeger–Gromoll splitting intuition** to an “asymptotic fiber-splitting” rigidity: a purely local Euclidean limit along the fibers would dictate the global geometry at infinity.  
   If false, constructing a counterexample would inform how nonnegative curvature can hide subtle global twisting even under asymptotically flat fiber models.  

Thus, the problem is mathematically sound, conceptually clear yet highly nontrivial, and sits at a natural frontier where geometric analysis, structure theory, and curvature rigidity meet.